% This version uses the latex2e styles, not the very ancient 2.09 stuff.
\documentclass[letterpaper,twocolumn,10pt]{article}
\usepackage{usenix,epsfig,endnotes,url}
\begin{document}

%don't want date printed
\date{}

%make title bold and 14 pt font (Latex default is non-bold, 16 pt)
\title{\Large \bf SIRIUS: A Revolutionary Approach to I/O and Storage for Exascale Scientific Computing}

%for single author (just remove % characters)
\author{
{\rm Your N.\ Here}\\
Your Institution
\and
{\rm Second Name}\\
Second Institution
% copy the following lines to add more authors
% \and
% {\rm Name}\\
%Name Institution
} % end author

%\author{
%\alignauthor {Jay Lofstead\footnotemark[1], Carlos Maltzahn\footnotemark[2], Scott Klasky\footnotemark[3], Hasan Abbasi\footnotemark[3], Qing Liu\footnotemark[3]\\
%Matthew Curry\footnotemark[1], Mark Ainsworth\footnotemark[5], Manish Parashar\footnotemark[4]}\\
%\footnotemark[1]~~Sandia National Laboratories
%\footnotemark[2]~~University of California\\
%\footnotemark[3]~~Oak Ridge National Laboratory
%\footnotemark[4]~~Rutgers University
%\footnotemark[5]~~Brown University\\
%}

\maketitle

% Use the following at camera-ready time to suppress page numbers.
% Comment it out when you first submit the paper for review.
\thispagestyle{empty}


\subsection*{Abstract}
The longest existing parallel I/O APIs, HDF5, PnetCDF, and NetCDF, are focused
on treating the entire output from a paralell application as a single unit.
These libraries work to gather the data from all processes and arrange it such
that it is organized much like as if it were written from a single process.
This approach, in particular the two-phase collective I/O with data sieving
portion, has proven difficult to scale for 3-D domain decompositions at scale.
ADIOS shifted the focus from treating the entire application output as a single
entity to treating the output from each writer as a self-contained entity. This
eliminated the need for the problematic two-phase collective I/O and has
demonstrated good scalability well into the petascale era. Another shift is
required to address the needs of exascale systems. New system I/O bandwidth is
not keeping pace with the compute acceleration leading to accute bottlenecks.
While ``burst buffer'' technology will help address the performance gap, it
does not address the problem completely. This paper describes a new approach
to not just an I/O library, but also the underlying storage infrastructure to
address the needs of exascacle applications cognizant of projected exascale
platform characteristics and in the context of current industry trends.

\section{Introduction}

High performance computing simulations have driven tightly coupled platform
scaling for decades. The output from these simulations are analyzed using
regular tools and techniques optimized for this environment.  The first
generation of standardized APIs for scientific computing IO focused on serial
output from a single process to a distributed or local file system.  These
systems still exist in the form of serial
HDF5~\cite{chilan:2006:hdf5-performance} and the end-of-life NetCDF
3.x~\cite{netcdf3}. The key feature of these APIs was the logical data model
that eliminated any parallel data decomoposition artifacts from the actual data
storage. This offered consistent performance and regular data access patterns.

The second generation is really an evolution of the first with HDF-5 adding
parallel capabilities and NetCDF bifurcating into
NetCDF4~\cite{rew:2011:netcdf-guide} and PnetCDF~\cite{torrellas:2003:pnetcdf}
with PnetCDF almost completely maintaining the data model and NetCDF4 shifting
to use HDF5 underneath. The other improvement is incorporating knowledge of an
underlying parallel file system to support large files written in unision from
many processes.

The third generation is a reaction to the second based on recognizing that
maintaining the logical data model in physical form became inordinately
expensive~\cite{lofstead:2011:six-degrees} and offered a performance penalty in
most cases when running at scale. Libraries like
ADIOS~\cite{lofstead:2009:adaptable} shifted from writing complete variables
from distributed processes as a single entity to treating all data from a
single source to be treated as a unit. This eliminated the data reorganization
at the cost of more complicated reading patterns. In spite of this extra
complication, the performance is generally
better~\cite{lofstead:2011:six-degrees}.

SIRIUS represents the next generation. New platforms are introducting
additional storage layers to accelerate IO operations and to compensate for
plateauing storage array performance. At the same time, data sizes continue
to grow. The fourth generation IO API will have to both address the IO API
and middleware layer as well as the deeper, more complex storage hierarchy.
SIRIUS considers the local portion of each globally distributed variable as a
management entity with data placement decisions across the entire storage
hierarhy rather than on a single tier.  Decisions on what data to place where
is driven by the data contents itself according to end-user specified
priorities and system characteristics. For example, if temperature and pressure
are nearly always retrieved together, organizing the data to keep just these
two arrays together without any interleaving from other variables or scalars
will accelerate later reading. If these variables could also be ``processed''
such that the important data features are maintained in high resolution while
less interesting regions are stored in highly compressed--even using lossy
compression--forms, the amount of data written can be greatly increased. While
accelerating writing is important. This approach also supports accelerating
reading.

This paper presents an overview of the efforts to build such a fourth
generation storage and IO system. A discussion of the end-user API, storage
challenges, quality of service features, and metadata challenges are all
discussed briefly.

%\begin{figure}[htbp]
%\vspace{-0.10in}
%\centering
%\includegraphics[width=\columnwidth]{images/arch-mapping}
%\vspace{-0.20in}
%\caption{Architecture and Component Mapping}
%\label{fig:arch-mapping}
%\vspace{-0.15in}
%\end{figure}

The rest of the paper is organized as follows. A brief overview of related work
is presented first in Section~\ref{sec:related}. Section~\ref{sec:features}
discusses the differentiating platform features SIRIUS will provide.
Section~\ref{sec:storage}
discusses the challenges and opportunities for deploying scalable storage
infrastructure.
Section~\ref{sec:qos}
discusses the quality of service features proposed for this project.
Next is an exploration of the metadata challenges required to support such a
system in Section~\ref{sec:metadata}.
An initial prototype demonstration of the functionality is presented in
Section~\ref{sec:evaluation}.
The paper is concluded in
Section~\ref{sec:conclusion} with a summary of the broad issues covered in the
paper.


\section{Related Work}
\label{sec:related}

Ceph~\cite{weil:ceph} is a distributed object store and file system. It offers
both a POSIX and object interface including features typically found in parallel
file systems.
Ceph's unique striping approach uses pseudo-random numbers with a
known seed eliminating the need for the metadata service to track where each
stripe in a parallel file is placed.
PVFS~\cite{carns:pvfs} offers optimizations to reduce metadata server load,
such as a single process opening a file and sharing the handle.
It has been
commercialized in recent years as OrangeFS.
Lustre~\cite{braam:lustre-arch} has become the de facto standard on most major
clusters offering scalable performance and fine-grained end-user and
programmatic control over how data is placed in the storage system.
GPFS~\cite{schmuck:gpfs} offers a hands-off approach for providing good
performance for scaling parallel IO tasks and is used extensively by its owner,
IBM.
Panasas~\cite{panasas:architecture} seeks to offer a dynamically adaptive
striping system that detects the need for additional stripes for performance
and adjusts the file layout as necessary.

Other file systems, like GoogleFS~\cite{ghemawat:googlefs} and
HDFS~\cite{Shvachko:2010:hdfs}, address distributed rather than parallel
computing and cannot be compared directly.
The primary difference between
distributed and parallel file systems is the ability of the file system to
store and retrieve data simultaneously from multiple clients, in parallel, and
treat the resulting collection of pieces as a single object.  Distributed file
systems rely on a single client creating a file, but distributing the set of
files across a wide array of storage devices.
The other, popular distributed
file system of note is NFS~\cite{powlowski:1994:nfs3} that has been used for
decades for enterprise file systems.

\section{Platform Features}
\label{sec:features}

The platform has two key features that drive the rest of the design. First is
the data refactoring/compression approach. Second are quality of service
estimates. 

\noindent\textbf{Refactoring} may be a uniform operation across an
entire variable. For example, splitting a double into three parts--two byes,
four bytes, and two bytes, yields a two-byte value that is the majority of the
data precision. The next four bytes would be most of the rest of the data
precision. The last two bytes are the remainder of the precision. For quick
overviews, just reading the two-byte portion is sufficient to get a rough idea
of data contents. If particular data looks interesting, more and more data can
be read to improve the precision. This can easily be applied uniformly by each
process without coordinating. From a storage hierarchy perspective, the primary
value portion could be written to high performance, low capacity storage while
the rest is written to slower storage. That would reduce the data quantity
written to NVM, for example, by 75\% for a 2/6 split enabling far more output
to be written to this fast tier with the rest spooled asynchronously.

A more complex, but interesting approach is to offer an {\em Auditor} that
performs a similar, but lower precision computation that greatly reduces
computation intensity with only a minor loss of precision over short spatial
and temporal domains. As the simulation progresses, the auditor drift would
need to be reset to bring it back into alignment. The much smaller, reduced
precision data could be used instead of the full precision in many
circumstances.  A more dynamic lossy approach would be to apply a wavelet
compression~\cite{klappenecker:1995:wavelet} approach that generates different
compressed data sizes at a fixed error rate depending on the data entropy.
These are still uniform operators, but the resulting data can be of radically
different sizes. For this approach, some or all of the lower precision data
can be written to the fast tier while full precision can be spooled
asynchrnously to the slower tier(s).

A middle ground approach between these two refactoring approaches is to use a
non-uniform operator across a variable. Instead of, for example, splitting on
byte boundaries to make multiple values, a function that identifies data
features through a local analytics function could identify ``interesting'' data
that should be maintained in full precision while other areas can be more
highly compressed. Such an operator would decompose the data from a process
further than a single process or node depending on the data features. This
approach may balance precision better than wavelet compression, but at a higher
analytics cost. We would put data areas with interesting features on the fast
tier while the rest could be spooled to the slow tier asynchronously.

We plan to investigate all three approaches as different applications will have
different techniques that work best. Our initial tests have demonstrated the
potential for the auditor approach in an isolated environment. We have also
built a prototype ADIOS transport method for the byte-split apporach. Any
performant storage system should offer the flexibility to incorporate
application specific data refactoring techniques and sufficient metadata
support for an arbitrary client to access the data successfully.

\noindent\textbf{Quality of Service} deadlines are important for driving how
much time is spent performing IO. At the same time, data quality must be
maintained at a sufficient level for the resulting science to be possible.
The system must offer a way for an end-user to negotiate with the system for
trading off between time spent in IO and the amount of data written or read.
Part of that tradeoff will be increased or reduced data quality, depending on
the system. For example, when data is written, a few different forms may be
output given sufficient space and time. One may be every fourth element
across a regular mesh reduces the data size to $1/64^{th}$ the full data set.
Further dividing that using byte-splitting can reduce it by another 75\%. For
quick overviews to search data sets, this reduced form can be sufficient
driving which data sets should be investigated in more detail. The key driving
factor is time. The applicaiton scientist can read 0.4\% of the full data
volume with a similar time reduction. SIRIUS will offer a negotiation process
to read any data.

First, a user will request a particular variable and timesetp. SIRIUS will
return a table similar to Table~\ref{tab:results}. For a given variable ``A'',
the end-user can then select the data quality and approximate reading time for
the operation. In this case, there are 5 options for obtaining the whole data
set (1/2/3/4, 5/6, 7/8, 9, 10). The first four represent a spatial and byte
split, 5/6 represent just strided writing, 7/8 represent just a byte-split, 9
represents a wavelet compressed version, and 10 is the full precision data.
The user would specify to read a particular data instance that would proceed
normally.  Quality of Service soft guarantees specified as part of the request
will result in SIRIUS allocating soft bandwidth reservations to meet the
request. While some of the error rates seem to render the data unusable,
consider a 3000x3000x3000 simulation domain. That would yield a 27 gigapixel
image should all of the data be used. Visualization will already radically
reduce the data sizes to make usable visualizations that work on a standard
display. A 99\% reduction would yield a 270 megapixel image--still quite large.

\begin{table}[tbp]
\centering
\caption{Reading Options for Variable ``A''}
\label{tab:results}
\begin{tabular}{|l|l|l|l|l|l|}
\hline
\# & Size & Time & Time Err & Refactoring & Data Err\\
\hline
1 & $(\frac{1}{4})^3(\frac{3}{8})$A & 10s & $\pm$ 3s & stride, byte-split & 99\% \\
\hline
2 & $(\frac{3}{4})^3(\frac{3}{8})$A & 90s & $\pm$ 30s & stride, byte-split & 58\% \\
\hline
3 & $(\frac{1}{4})^3(\frac{5}{8})$A & 16s & $\pm$ 5s & stride, byte-split & 0.01\% \\
\hline
4 & $(\frac{3}{4})^3(\frac{5}{8})$A & 120s & $\pm$ 50s & stride, byte-split & 0.01\% \\
\hline
\hline
5 & $(\frac{1}{4})^3$A & 1200s & $\pm$ 30s & stride & 98\% \\
\hline
6 & $(\frac{3}{4})^3$A & 2400s & $\pm$ 90s & stride & 58\% \\
\hline
\hline
7 & $(\frac{3}{8})$A & 1350s & $\pm$ 120s & byte-split & 5\% \\
\hline
8 & $(\frac{5}{8})$A & 2250s & $\pm$ 120s & byte-split & 0.01\% \\
\hline
\hline
9 & A & 36s & $\pm$ 6s & wavelet & 1\% \\
\hline
\hline
10 & A & 3600s & $\pm$ 600s & none & 0\% \\
\hline
\end{tabular}
\end{table}


\section{Storage Challenges and Opportunities}
\label{sec:storage}

The supercomputing platforms being deployed today are all incorporating a
Non-volatile memory (NVM) layer as a fast cache between compute and the storage
array. The assumption is that nearly all IO will be performed against this
layer to hide the much greater latency to the storage array. The challenge is
the NVM capacity is expensive limiting the amount deployed. The unspoken
limitation is that this fast layer still must drain to the relatively slow
storage array. This performance mismatch will limit capacity and performance
just due to the drain and data stage-in operations.

A more effective way to use these resources is to only store the most
important data in the fast tier and store lower precision representations of
the rest of the data in the slower tiers. This lowers pressures on the NVM
layer increasing performance overall by greatly reducing the drain/stage-in
operations and reducing the data volumes stored on a per application basis.
As part of the data refactoring, SIRIUS will use the entire storage hierarchy
to store different portions of each variable as illustrated in
Table~\ref{tab:results}. For example, for the 1/2/3/4 setup, 4 different tiers
may be used to store this single variable.

\section{Quality of Service Features}
\label{sec:qos}

SIRIUS offers a novel feature related to managing IO time as a first class
characteristic. Rather than being at the mercy of the storage array
performance, an end-user can describe the local data to the storage system and
receive an estimate on how long that data will take to write. If that is
acceptable within a margin of error, the API can write. If it is unacceptable,
the end-user can make decisions about skipping the output or maybe spending a
little time with a strong data refactoring to reduce the data size.

On reading, a user may specify that they are willing to wait for an approximate
time for an operation. This will negotiate with the storage system to determine
the data quality returned. If the requested data quality cannot be retrieved
within a given time, the storage array may be able to apply an additional
data reduction operator to decrease the amount of time spent in IO.
Alternatively, if multiple different refactored data versions are stored,
alternative versions can be requested instead.

\section{Metadata Challenges}
\label{sec:metadata}

The metadata system has a daunting challenge. Traditional metadata systems for
both distributed and parallel file systems only have to deal with a single name
space and storage hierarchy. While hierarichal storage management~\cite{blaze:1992:hsm}
systems offer a way to show a single name space, the reality is that the tier
closest to the user is a cache for the other tiers with data migrating as
necessary to support user interaction. SIRIUS is working in a fundamentally
different way. Each tier is treated as a first class citizen with direct
client access possible. Where data is actually retrieved from is where it is
actually stored rather than forcing data migration and caching. Systems like
Sirocco~\cite{curry:2015:sirocco} offer a way to support this model directly and are a
fundamental part of our system.

The additional challege SIRIUS metadata services face is incorporating
user-specified attributes including data refactoring functions. Since an
arbitrary user must be able to access any data written for which they have
permissions. Having data locked behind data compression that has changed over
time may result in effectively encrypting data and losing the decryption key.

We also seek to offer the ability to also store multiple versions of the same
data using different refactoring techniques. This will offer additional options
for higher precision OR higher performance reading.

\section{Demonstration}
\label{sec:evaluation}

This stack has an early prototype implementation intended to test concepts
rather than performance and scalability. It has focused on examining the
interaction of the different APIs for each layer to flesh out any detailed
requirements or concerns that may have been missed in the conceptualization of
this IO stack. To demonstrate the viability of the IO stack described in this
paper, we show some very early performance results from the untuned prototype.

\section{Conclusions}
\label{sec:conclusion}

Overall, the SIRIUS storage system is ushering in a new generation of both IO
APIs and storage system managed together. It incorporates additional metadata
to help accelerate data access and ultimately time to insight. Our initial
efforts show promise that such a system is both possible and can offer
attractive performance to better use the supercomputing platform.

%\subsection{How to Build Your Paper}
%
%You have to run {\tt latex} once to prepare your references for
%munging.  Then run {\tt bibtex} to build your bibliography metadata.
%Then run {\tt latex} twice to ensure all references have been resolved.
%If your source file is called {\tt usenixTemplate.tex} and your {\tt
%  bibtex} file is called {\tt usenixTemplate.bib}, here's what you do:
%{\tt \small
%\begin{verbatim}
%latex usenixTemplate
%bibtex usenixTemplate
%latex usenixTemplate
%latex usenixTemplate
%\end{verbatim}
%}
%

\section{Acknowledgments}

%\includegraphics[scale=0.07]{logos/doe_logo}
%\includegraphics[scale=0.30]{logos/snl_logo}
%\includegraphics[scale=0.35]{logos/nnsa_logo}
Sandia National Laboratories is a multi-program laboratory managed and operated
by Sandia Corporation, a wholly owned subsidiary of Lockheed Martin
Corporation, for the U.S. Department of Energy's National Nuclear Security
Administration under contract DE-AC04-94AL85000.

{\footnotesize \bibliographystyle{acm}
\bibliography{paper}}

%\theendnotes

\end{document}
